\documentclass[12pt]{article}
\usepackage[T1]{fontenc}
\usepackage[utf8]{inputenc}
\usepackage{lmodern}
\usepackage{textcomp}
\usepackage{enumitem}
\usepackage{geometry}
\usepackage{graphicx}
\usepackage{amsmath, amssymb}
\geometry{margin=1in}
\begin{document}
\begin{center}
Economics - Graduate Level Assessment 22 \
Total Marks: 40 \
Answer all questions.
\end{center}

\section*{Multiple Choice (10 marks)}
\begin{enumerate}[label=\arabic*.]
\item If nominal GDP equals $5000 and real GDP equals $4000 then the GDP Deflator equals
\begin{enumerate}[label=(\alph*)]
    \item 0.5
    \item 125
    \item 500
    \item 0.8
    \item 100
\end{enumerate}
\item What is the effect on savings of a tax cut of \$10 billion? Is thisinflationary or deflationary? Assume that the marginal propensity to consume is 0.85.
\begin{enumerate}[label=(\alph*)]
    \item \$1.5 billion
    \item \$8.5 billion
    \item \$0.85 billion
    \item \$15 billion
    \item \$56.67 billion
\end{enumerate}
\item The real interest rate is
\begin{enumerate}[label=(\alph*)]
    \item the anticipated inflation rate plus the nominal interest rate
    \item the nominal interest rate multiplied by the inflation rate
    \item what one sees when looking at bank literature
    \item the anticipated inflation rate minus the nominal interest rate
    \item the nominal interest rate minus anticipated inflation
\end{enumerate}
\item What do we mean when we say that the farm problem may be cor-rectly envisioned as a problem of resource misallocation?
\begin{enumerate}[label=(\alph*)]
    \item Overproduction and wastage of crops
    \item Lack of alternative employment
    \item Excessive government intervention in agriculture
    \item Inefficient farming methods
    \item Low income of farmers
\end{enumerate}
\item What will be the properties of the OLS estimator in the presence of multicollinearity?
\begin{enumerate}[label=(\alph*)]
    \item It will not be unbiased or efficient
    \item It will be consistent and efficient but not unbiased
    \item It will be efficient but not consistent or unbiased
    \item It will be efficient and unbiased but not consistent
    \item It will be consistent and unbiased but not efficient
\end{enumerate}
\end{enumerate}

\section*{True / False (10 marks)}
\textit{Answer True or False}
\begin{enumerate}[label=\arabic*.]
\setcounter{enumi}{5}
\item True or False: A price decrease for a normal good will always lead to an increase in quantity demanded, regardless of the income effect.
\item True or False: An increase in the price of tea, a substitute for coffee, will likely lead to an increase in the demand for coffee.
\item True or False: The removal of a protective tariff on imported steel leads to improved allocative efficiency in the steel market.
\item True or False: GDP is unaffected by the purchase of a second-hand painting, as it does not contribute to current economic production.
\item True or False: The entry of new firms into a monopolistically competitive industry will not affect the overall costs of production in the market.
\end{enumerate}

\section*{Long Form Response (20 marks)}
\textit{Answer each question with a clear and complete explanation.}
\begin{enumerate}[label=\arabic*.]
\setcounter{enumi}{10}
\item Discuss the concept of structural unemployment in the context of technological advancements, using the example of a dark room technician losing their job due to declining film camera usage.
\item Discuss the rationale behind subtracting expenditures for imports when measuring aggregate demand, and analyze how this differs from the inclusion of exports in the calculation.
\end{enumerate}

\vfill
\noindent\textit{End of Paper}
\end{document}